\section{Problem 7: Tight Binding}

In Problems 5 and 6, you studied free-electron models of metals. The Drude model treated electrons as classical carriers, and the Sommerfeld model added Fermi-Dirac statistics to free electrons. In this problem, you will replace the free-electron spectrum with an electronic band generated from coupled atomic orbitals.

The tight-binding model shows how electronic bands arise from coupling localized atomic orbitals across a periodic lattice. A finite chain gives discrete molecular-like energy levels. A periodic infinite chain gives a continuous band \(E(k)\). Increasing the number of sites connects the finite matrix spectrum to the band-structure picture.

The computational chain is
\[
\epsilon_0,t,a
\longrightarrow
H
\longrightarrow
E_j
\longrightarrow
E(k)
\longrightarrow
D(E)
\longrightarrow
\text{band filling}.
\]

Here, \(\epsilon_0\) is the onsite orbital energy, \(t\) is the nearest-neighbor hopping parameter, \(a\) is the lattice spacing, \(H\) is the tight-binding Hamiltonian matrix, \(E_j\) are finite-chain eigenvalues, \(E(k)\) is the analytical band dispersion, and \(D(E)\) is the density of states.

\subsection{Physical model}

Consider a one-dimensional chain of \(N\) identical atoms. Each atom contributes one localized orbital \(|n\rangle\), where \(n=0,1,\ldots,N-1\).

Assume that the localized orbitals are orthonormal:
\[
\langle i|j\rangle=\delta_{ij}.
\]
With this assumption, the tight-binding eigenvalue problem has the ordinary matrix form
\[
H\mathbf c=E\mathbf c.
\]
You do not need to use an overlap matrix in this problem.

The onsite energy of each orbital is \(\epsilon_0\). Neighboring orbitals are coupled by a hopping parameter \(t\). Use the sign convention
\[
H_{i,i}=\epsilon_0,
\qquad
H_{i,i+1}=H_{i+1,i}=-t.
\]
For \(t>0\), this convention gives a band minimum at \(k=0\).

The nearest-neighbor tight-binding Hamiltonian is
\[
\hat H
=
\epsilon_0\sum_n |n\rangle\langle n|
-
t\sum_n
\left(
|n\rangle\langle n+1|
+
|n+1\rangle\langle n|
\right).
\]

In the localized orbital basis, the Hamiltonian matrix elements are
\[
H_{ij}
=
\epsilon_0\delta_{ij}
-
t
\left(
\delta_{i,j+1}
+
\delta_{i,j-1}
\right),
\]
with the boundary terms determined by the chosen boundary condition.

\subsection{Boundary conditions}

Implement both open and periodic boundary conditions.

For open boundary conditions, the first and last sites are not connected. The Hamiltonian has hopping terms only between neighboring sites inside the finite chain.

For periodic boundary conditions, the last site is connected back to the first site. This represents a ring and approximates an infinite crystal more directly:
\[
H_{0,N-1}=H_{N-1,0}=-t.
\]

The analytical band dispersion in this problem is the periodic infinite-chain benchmark. Open chains have standing-wave-like eigenstates and boundary-modified finite-size eigenvalues. Periodic chains have eigenvalues that directly sample the Bloch band.

\subsection{Analytical band dispersion}

For an infinite periodic chain, use the Bloch form
\[
\psi_n(k)=e^{ikna}.
\]

Substituting this into the tight-binding Hamiltonian gives
\[
E(k)=\epsilon_0-2t\cos(ka),
\qquad
-\frac{\pi}{a}\le k\le\frac{\pi}{a}.
\]

The band center is \(\epsilon_0\). For \(t>0\), the band minimum occurs at \(k=0\), and the band maximum occurs at \(k=\pm\pi/a\).

The bandwidth is
\[
W=E_{\max}-E_{\min}=4|t|.
\]

For a finite periodic chain, the allowed wavevectors are discrete:
\[
k_m=\frac{2\pi m}{Na}.
\]
These \(k_m\) values should be mapped into the first Brillouin zone. The periodic finite-chain eigenvalues should agree with the analytical dispersion evaluated at these allowed wavevectors.

\subsection{Connection to the free-electron model}

The Sommerfeld model used a free-electron spectrum. In one dimension, the free-electron energy has the form
\[
E_{\mathrm{free}}(k)=\frac{\hbar^2k^2}{2m}.
\]

The tight-binding model replaces this with a lattice-generated band,
\[
E_{\mathrm{TB}}(k)=\epsilon_0-2t\cos(ka).
\]

This changes the density of states, the bandwidth, and the way the Fermi energy moves as the band is filled. Problem 7 therefore connects the electronic statistics from Problem 6 to a band structure generated by a lattice Hamiltonian.

\subsection{Numerical implementation}

Complete the provided skeleton file \texttt{tight\_binding.py} in \texttt{./src/dlsu\_solst}. Implement a class named \texttt{TightBinding1D}.

The constructor should store the number of sites \(N\), onsite energy \(\epsilon_0\), hopping parameter \(t\), lattice spacing \(a\), and boundary condition.

The method \texttt{hamiltonian\_matrix} should construct the finite-chain Hamiltonian matrix \(H\).

The method \texttt{eigenvalues} should diagonalize \(H\) and return the finite-chain energy eigenvalues sorted from lowest to highest.

The method \texttt{k\_grid} should return a grid of wavevectors over the first Brillouin zone:
\[
-\frac{\pi}{a}\le k\le\frac{\pi}{a}.
\]

The method \texttt{allowed\_k\_values} should return the discrete \(k_m\) values for a finite periodic chain:
\[
k_m=\frac{2\pi m}{Na},
\]
mapped into the first Brillouin zone.

The method \texttt{analytical\_dispersion(k)} should compute
\[
E(k)=\epsilon_0-2t\cos(ka).
\]

The method \texttt{density\_of\_states} should estimate the spinless density of states from the finite-chain eigenvalues. The density of states should be normalized so that
\[
\int D(E)\,dE\approx N.
\]

The method \texttt{fermi\_energy(n\_electrons)} should estimate the Fermi energy by filling the finite-chain eigenvalues from lowest to highest. The Hamiltonian itself is spinless, so it gives \(N\) spatial-orbital energy levels. Include spin degeneracy only when counting electrons for band filling. Therefore, each spatial level can hold two electrons.

For this problem, define the finite-chain Fermi energy as the highest occupied eigenvalue.

The implementation should be progressive. The eigenvalues should use the Hamiltonian matrix. The density of states should use the eigenvalues. The Fermi energy should use the sorted eigenvalues and the electron count.

\subsection{Numerical checks}

Check that the Hamiltonian matrix is Hermitian:
\[
H=H^\dagger.
\]

Check that all numerical eigenvalues are real.

For periodic boundary conditions, check that the numerical eigenvalues agree with the analytical dispersion evaluated at the allowed finite-chain wavevectors:
\[
E_j^{\mathrm{num}}
\approx
E(k_m).
\]

Check that the analytical bandwidth satisfies
\[
W\approx4|t|.
\]

Check that changing \(\epsilon_0\) shifts the entire band without changing the bandwidth.

Check that increasing \(|t|\) increases the bandwidth.

Check that the spinless density of states integrates to the expected number of spatial states:
\[
\int D(E)\,dE\approx N.
\]

Check that the band can hold \(2N\) electrons when spin degeneracy is included.

\subsection{Numerical experiments}

First, compute the Hamiltonian matrix and eigenvalues for a finite chain. Compare open and periodic boundary conditions for the same \(N\), \(\epsilon_0\), \(t\), and \(a\).

Second, compute the analytical dispersion \(E(k)\) over the first Brillouin zone. Overlay the periodic finite-chain eigenvalues at the allowed \(k_m\) values. This should show that the periodic finite-chain spectrum samples the analytical band.

Third, study the effect of system size \(N\). Increase \(N\) and show that the finite-chain eigenvalues become a denser sampling of the continuous band.

Fourth, study the effect of the hopping parameter \(t\). Compute the band for several values of \(t\) and show that the bandwidth scales as \(4|t|\).

Fifth, estimate the density of states from the finite-chain eigenvalues. Plot the spinless density of states for several values of \(N\). The one-dimensional tight-binding density of states is not flat; it has strong structure near the band edges. A finite-chain histogram will approximate this structure depending on the number of sites and the bin width.

Finally, study band filling. For electron counts from \(0\) to \(2N\), fill the finite-chain eigenvalues from lowest to highest using spin degeneracy. Estimate the Fermi energy as the highest occupied eigenvalue. This shows how the Fermi level rises as the band is filled.

\subsection{Required plots}

Include the following plots:

\begin{enumerate}
    \item analytical dispersion \(E(k)\) over the first Brillouin zone;
    \item periodic finite-chain eigenvalues overlaid on the analytical dispersion at the allowed \(k_m\) values;
    \item open-boundary and periodic-boundary finite-chain eigenvalues;
    \item band dispersion for several values of \(t\);
    \item bandwidth \(W\) versus \(|t|\);
    \item spinless density of states \(D(E)\) for several chain lengths \(N\);
    \item Fermi energy versus electron count.
\end{enumerate}

Each plot should have labeled axes, units where appropriate, and a short caption or explanation.

\subsection{Results}

Your results should report the Hamiltonian matrix structure, the finite-chain eigenvalues, and the analytical dispersion relation.

Report the numerical checks. State whether \(H\) is Hermitian, whether the eigenvalues are real, whether the periodic-chain eigenvalues agree with the analytical dispersion at the allowed \(k_m\) values, and whether the bandwidth is approximately \(4|t|\).

Report how the spectrum changes when the boundary condition changes from open to periodic.

Report how the spectrum changes as \(N\) increases. Explain why the finite-chain eigenvalues become a denser sampling of the band.

Report how changing \(\epsilon_0\) shifts the band and how changing \(t\) changes the bandwidth.

Report the spinless density of states and state whether it integrates to \(N\).

Report the Fermi energy as a function of electron count and explain how the Fermi level moves through the band as the band is filled.

\subsection{Discussion}

In your discussion, explain how tight binding differs from the free-electron model. The free-electron model starts from delocalized plane waves, while tight binding starts from localized atomic orbitals and builds bands by coupling neighboring orbitals.

Explain the physical meaning of the onsite energy \(\epsilon_0\). It sets the center of the band.

Explain the physical meaning of the hopping parameter \(t\). It controls how strongly neighboring orbitals are coupled. Larger \(|t|\) produces a wider band.

Explain why a finite chain has a discrete set of eigenvalues, while an infinite periodic chain has a continuous band \(E(k)\).

Explain why periodic boundary conditions are useful for modeling bulk band structure.

Discuss the spin convention. The Hamiltonian is constructed for one spin channel. Spin degeneracy is included only when counting electrons for band filling.

Discuss band filling. At low electron count, electrons occupy the bottom of the band. As the electron count increases, the Fermi energy rises through the band. A completely filled spin-degenerate band contains \(2N\) electrons for \(N\) spatial orbitals.

Conclude by summarizing the computational chain
\[
\epsilon_0,t,a
\longrightarrow
H
\longrightarrow
E_j
\longrightarrow
E(k)
\longrightarrow
D(E)
\longrightarrow
\text{band filling}.
\]
This chain shows how coupled atomic orbitals generate electronic bands and how band filling connects the tight-binding spectrum to the Fermi energy.

\subsection{Collaboration reminder}

No points will be deducted for appropriate collaboration, provided that the collaboration is disclosed and the submitted code, calculations, plots, explanations, and conclusions are your own.